% !TeX root = ../../Book3.tex
\section{基本推论}
\label{b3:sec:0204}

Cauchy 估计中可以选择的半径，能把局部导数变成全局增长的探测工具。对整函数把圆周不断扩大，增长过慢会迫使高阶导数消失；在有界区域内，把圆周平均与模的上界比较，又会得到内部极值的限制。本节沿这两个方向证明 Liouville 定理、代数基本定理与最大模原理，最后在单位圆盘中加入中心归一化，得到 Schwarz 引理。一般两点之间的距离估计留待\chapref{b3:ch:06}。

\subsection{Liouville 定理}
\label{b3:subsec:020401}

实轴上的有界光滑函数有很多，例如正弦函数。若同时要求在整个复平面全纯，有界性却会使所有变化消失。\term[liouville-dingli]{Liouville 定理}[Liouville's theorem]的关键是围道可以取得任意大，而函数的模始终受同一个常数控制。

\Needspace{6\baselineskip}
\begin{theorem}[Liouville]\label{b3:thm:02-liouville}
设 \(f:\mathbb C\to\mathbb C\) 为整函数。若存在常数 \(M\ge0\)，使 \(|f(z)|\le M\) 对所有 \(z\in\mathbb C\) 成立，则 \(f\) 为常函数。
\end{theorem}
\begin{proof}
固定任意 \(a\in\mathbb C\)。因为 \(f\) 在整个平面全纯，对每个 \(r>0\) 都能在圆周 \(|z-a|=r\) 上使用\cref{b3:thm:02-cauchy-estimate}，得到
\[
 |f'(a)|\le\frac{M}{r}.
\]
左端与半径无关，令 \(r\to\infty\) 得 \(f'(a)=0\)。\(a\) 任意，故 \(f'\) 在整个平面为零。对任意两点 \(p,q\)，沿连接它们的线段使用\cref{b3:prop:02-primitive-endpoints}，有
\[
 f(q)-f(p)=\int_{[p,q]}f'(z)\,\mathrm{d}z=0.
\]
于是 \(f\) 为常数。
\end{proof}

整性与全平面的共同上界各自不可缺少。函数 \(z\) 在单位圆盘有界且全纯，却不能在该圆盘内把取样半径无限扩大。整函数 \(e^{\mathrm{i}z}\) 在实轴上模恒为一，但 \(|e^{\mathrm{i}(x+\mathrm{i}y)}|=e^{-y}\)，沿负虚轴无界；一条直线上的界不能控制大圆周。Liouville 定理限制的是函数在全平面的同时增长。

\ProseParagraphBreak

同一证明还说明，允许多项式速度的增长，只会留下有限多个 Taylor 系数。这里要固定求导阶数，再把半径放大；不同阶数允许分别取极限。

\Needspace{6\baselineskip}
\begin{proposition}[多项式增长判别]\label{b3:prop:02-polynomial-growth}
设 \(f:\mathbb C\to\mathbb C\) 为整函数，\(m\ge0\) 为整数。若存在常数 \(C\ge0\)，使
\[
 |f(z)|\le C(1+|z|)^m\quad(z\in\mathbb C),
\]
则 \(f\) 是次数不超过 \(m\) 的多项式，允许为零多项式。反之，每个次数不超过 \(m\) 的多项式都满足这样的估计。
\end{proposition}
\begin{proof}
固定整数 \(n>m\)。在以零为中心、半径为 \(r\) 的圆周上，Cauchy 估计给
\[
 \frac{|f^{(n)}(0)|}{n!}
 \le\frac{C(1+r)^m}{r^n}
 =C\left(1+\frac1r\right)^m r^{m-n}\longrightarrow0
 \quad(r\to\infty).
\]
所以所有阶数大于 \(m\) 的 Taylor 系数为零。\(f\) 的 Taylor 展开在每个有限圆盘内都等于 \(f\)，而这些圆盘覆盖整个平面，故
\[
 f(z)=\sum_{n=0}^m\frac{f^{(n)}(0)}{n!}z^n
 \quad(z\in\mathbb C).
\]
反之，若 \(f(z)=\sum_{n=0}^m a_nz^n\)，则 \(|z|^n\le(1+|z|)^m\) 对 \(0\le n\le m\) 成立，从而取 \(C=\sum_{n=0}^m|a_n|\) 即得所需上界。
\end{proof}

这一判别把增长条件变成有限个未知系数。例如设整函数 \(f\) 满足 \(|f(z)|\le C(1+|z|)^2\)，且 \(f(0)=1\)、\(f(1)=f(-1)=2\)。先由命题写成 \(f(z)=a+bz+cz^2\)，再代入三点，得到
\[
 a=1,\quad b+c=1,\quad -b+c=1.
\]
故 \(f(z)=1+z^2\)。若没有增长条件，有限个点值无法唯一决定整函数；在这个解上加 \(z(z^2-1)e^z\)，三个指定点值仍相同。这里真正减少自由度的是全平面的增长控制。

\subsection{代数基本定理}
\label{b3:subsec:020402}

要证明一个多项式具有零点，可以反过来考察假设无零时的倒数。多项式在无穷远处变大，使倒数在外部有界；无零假设则使倒数在剩余紧集上有界。两部分合在一起，便能使用 Liouville 定理。这给出\term[daishu-jiben-dingli]{代数基本定理}[fundamental theorem of algebra]的分析证明。

\Needspace{6\baselineskip}
\begin{theorem}[代数基本定理]\label{b3:thm:02-fundamental-algebra}
设 \(n\ge1\) 为整数，\(P(z)=a_nz^n+\cdots+a_0\) 为复系数多项式，且 \(a_n\ne0\)。则存在复数 \(\alpha_1,\ldots,\alpha_n\)，允许重复，使
\[
 P(z)=a_n\prod_{j=1}^n(z-\alpha_j).
\]
因此 \(P\) 在复平面中计重数恰有 \(n\) 个零点。
\end{theorem}
\begin{proof}
先证明至少有一个零点。反设 \(P(z)\ne0\) 于整个平面，则 \(1/P\) 是整函数。对 \(z\ne0\)，写成
\[
 P(z)=a_nz^n\left(1+\sum_{k=0}^{n-1}\frac{a_k}{a_n}z^{k-n}\right).
\]
括号内有限和的每一项都在 \(|z|\to\infty\) 时趋于零。故可选 \(R>0\)，使
\[
 \sum_{k=0}^{n-1}\left|\frac{a_k}{a_n}\right|\cdot|z|^{k-n}
 \le\frac12\quad(|z|\ge R).
\]
于是外圆之外有
\[
 |P(z)|\ge\frac{|a_n|}{2}\cdot|z|^n,
 \quad\left|\frac1{P(z)}\right|
 \le\frac{2}{|a_n|\cdot R^n}.
\]
在闭圆盘 \(\overline{D(0,R)}\) 上，\(|P|\) 是连续的正函数；紧性使它取得严格正的最小值 \(\delta\)。因而这个圆盘上 \(|1/P|\le1/\delta\)。内外两界合起来说明 \(1/P\) 在整个平面有界。由 Liouville 定理，它为常数；该常数不为零，故 \(P\) 也为常数，与 \(n\ge1\) 矛盾。

既然已有零点 \(\alpha_1\)，多项式除法给出
\(P(z)=(z-\alpha_1)Q(z)\)，其中 \(Q\) 的次数为 \(n-1\)，最高次系数仍为 \(a_n\)。这一步也可逐项使用 \(z^k-\alpha_1^k\) 被 \(z-\alpha_1\) 整除的恒等式，因为余项 \(P(\alpha_1)\) 为零。对 \(Q\) 重复上述存在性结论，有限次之后得到所列的 \(n\) 个一次因子。

最后核对重数。若某个值 \(\alpha\) 在这些因子中出现 \(l\) 次，将它们合并可写成 \(P(z)=(z-\alpha)^lQ_\alpha(z)\)，且 \(Q_\alpha(\alpha)\ne0\)。由\cref{b3:prop:01-zero-factorization}，\(l\) 正是全纯零点的阶数。乘积表示也表明不存在未列出的零点，所以计重数总共为 \(n\)。
\end{proof}

例如 \((z-1)^2(z+2)\) 只有两个不同零点，但在 \(1\) 与 \(-2\) 的阶数分别为二和一，总重数仍为三。定理中的计数不要求各根不同，也不保证它们能用根式写出。

\ProseParagraphBreak

对任意固定 \(w\in\mathbb C\)，多项式 \(P-w\) 仍具有相同的最高次项。定理因此还说明：非常值复多项式把平面满射到平面，对每个值恰有 \(n\) 个计重数的原像。这个结论包含一项全局信息，不能只从局部开映射推出。例如指数函数在每点附近都可逆，其像也开放，却始终遗漏零；多项式的无穷远增长在这里排除了遗漏值。

\subsection{最大模原理}
\label{b3:subsec:020403}

在有界区域内，通常不能把半径推向无穷；可以利用的是圆周平均。全纯函数在中心等于周围值的平均，如果中心的模已经最大，平均中的各个值就没有偏离它的余地。\term[zuidamo-yuanli]{最大模原理}[maximum modulus principle]由此把内部极值与边界控制联系起来。

\Needspace{6\baselineskip}
\begin{theorem}[最大模原理]\label{b3:thm:02-maximum-modulus}
设 \(U\subseteq\mathbb C\) 为区域，\(f:U\to\mathbb C\) 全纯。若 \(|f|\) 在 \(U\) 内某点取得局部最大值，则 \(f\) 为常数。此外，设 \(V\subseteq\mathbb C\) 为非空有界区域，\(F:\overline V\to\mathbb C\) 连续且在 \(V\) 上全纯，则
\[
 \max_{z\in\overline V}|F(z)|
 =\max_{z\in\partial V}|F(z)|.
\]
\end{theorem}
\begin{proof}
设 \(a\in U\) 是 \(|f|\) 的局部最大点，记 \(M=|f(a)|\)。取 \(\varepsilon>0\)，使 \(D(a,\varepsilon)\subseteq U\)，且该圆盘上 \(|f|\le M\)。若 \(M=0\)，则 \(f\) 已在这个圆盘恒为零，由恒等定理在 \(U\) 上也恒为零。

若 \(M>0\)，取单位复数 \(\lambda=\overline{f(a)}/M\)，令 \(g=\lambda f\)。这样 \(g(a)=M\)，并且 \(|g|\le M\) 于 \(D(a,\varepsilon)\)。对每个 \(0<r<\varepsilon\)，圆盘内的 Cauchy 公式给出中心平均
\[
 M=\frac1{2\pi}\int_0^{2\pi}g(a+re^{\mathrm{i}t})\,\mathrm{d}t.
\]
取实部后，连续非负函数 \(M-\operatorname{Re}g(a+re^{\mathrm{i}t})\) 的积分为零，所以它处处为零。圆周上 \(\operatorname{Re}g=M\)，同时 \(|g|\le M\)，从
\(M^2+(\operatorname{Im}g)^2=|g|^2\le M^2\)
可知虚部也为零。因此 \(g=M\) 于这条圆周。让 \(r\) 取遍 \((0,\varepsilon)\)，得到 \(f\) 在整个小圆盘上为常数，再由恒等定理将常值延至 \(U\)。

对于第二项，\(\overline V\) 紧，连续函数 \(|F|\) 在其中取得最大值。如果最大点位于边界，等式立即成立；如果最大点位于内部，第一项使 \(F\) 在 \(V\) 为常数。\(V\) 非空且有界，其边界非空，连续性把同一个常值延续到边界，等式仍成立。
\end{proof}

第一章的开映射定理也解释了这一结论：非常值全纯函数把一个内点邻域映成开集，而任何包含 \(f(a)\) 的小圆盘都有模比 \(|f(a)|\) 更大的点。上面的平均值证明进一步指出取等为何迫使整个小圆周上的函数值相同。

\ProseParagraphBreak

边界最大模结论不要求区域单连通，也不要求边界光滑。例如在圆环 \(V=\{1<|z|<2\}\) 上，\(F(z)=z+1/z\) 全纯并连续到闭包。内圆周上 \(|F|\le2\)，外圆周上 \(|F|\le2+1/2=5/2\)，所以整个闭环上都有 \(|F|\le5/2\)。在 \(z=2\) 时达到 \(5/2\)，故这正是最大值。两条边界都要比较；最大模原理本身不替我们决定哪一条边界控制结果。

\ProseParagraphBreak

区域有界这一条件不能直接删去。在上半平面，函数 \(F(z)=e^{-\mathrm{i}z}\) 连续达到实轴，在实轴上模恒为一，但对 \(z=x+\mathrm{i}y\)、\(y>0\)，有 \(|F(z)|=e^y>1\)。它的增长出现在无穷远处，有限边界上的上界不足以控制内部。另一方面，对有界区域也须具有陈述中的连续边界值，才能在紧闭包上取到最大值。

\ProseParagraphBreak

最大模原理给出边界值的唯一性与稳定性。设 \(V\) 为有界区域，\(F,G\) 在 \(\overline V\) 上连续、在 \(V\) 上全纯，且边界上 \(|F-G|\le\eta\)，其中 \(\eta\ge0\)。将定理用于 \(F-G\)，便有
\[
 \max_{\overline V}|F-G|\le\eta.
\]
特别地，相同的连续边界值至多对应一个全纯函数。然而连续边界数据未必都有全纯解。若单位圆上的数据 \(\overline z\) 是某个 \(h\in C(\overline{\mathbb D})\)、在 \(\mathbb D\) 全纯的函数的边界值，则 \(zh(z)-1\) 在圆周恒为零。唯一性迫使 \(zh(z)-1\) 在圆盘恒零，但代入 \(z=0\) 得到 \(-1=0\)，矛盾。

\ProseParagraphBreak

内部极值还限制插值。例如 \(f\in C(\overline{\mathbb D})\) 在 \(\mathbb D\) 全纯，且圆周上 \(|f|\le1\)。若再指定 \(f(0)=1\)，则 \(f\equiv1\)，不能同时要求 \(f(1/2)=0\)。若只指定 \(f(0)=1/2\)，则没有这项刚性；函数 \((1+z)/2\) 就满足条件且非常值。起作用的是内点已经取到允许的最大模。

\ProseParagraphBreak

最小模的情形需要先排除所考察点为零点。设 \(f\) 在区域 \(U\) 全纯，\(|f|\) 于 \(a\in U\) 取得局部最小值且 \(f(a)\ne0\)。由连续性，\(f\) 在 \(a\) 的一个小圆盘上无零，因此 \(1/f\) 在该圆盘全纯，并在 \(a\) 取得局部最大模。最大模原理使 \(1/f\)、从而 \(f\) 在该圆盘常值，再由恒等定理得 \(f\) 在 \(U\) 常值。若允许 \(f(a)=0\)，函数 \(f(z)=z\) 就在原点取得真正的最小模而不恒定。这项条件仅在使用倒数的邻域内需要，不能用不存在的 \(1/f(a)\) 代替论证。

\subsection{Schwarz 引理初步}
\label{b3:subsec:020404}

对圆盘中的全纯函数，边界上界控制内部模；若还知道中心是一个零点，就应当能把这个零点带来的衰减提出来。\term[schwarz-yinli]{Schwarz 引理}[Schwarz lemma]正是对商 \(f(z)/z\) 使用最大模原理。困难只在于函数未必连续延至单位圆周，因此先在内部半径上估计，再逐渐扩大半径。

\Needspace{6\baselineskip}
\begin{theorem}[Schwarz]\label{b3:thm:02-schwarz}
记 \(\mathbb D=D(0,1)\)。若 \(f:\mathbb D\to\mathbb D\) 全纯且 \(f(0)=0\)，则
\[
 |f(z)|\le|z|\quad(z\in\mathbb D),
 \quad |f'(0)|\le1.
\]
若存在 \(z_0\in\mathbb D\setminus\{0\}\) 使 \(|f(z_0)|=|z_0|\)，或 \(|f'(0)|=1\)，则存在常数 \(c\in\mathbb C\)、\(|c|=1\)，使 \(f(z)=cz\) 于整个 \(\mathbb D\)。反之，每个这样的旋转都取等。
\end{theorem}
\begin{proof}
因为 \(f(0)=0\)，其 Taylor 展开在原点没有常数项，所以
\[
 g(z)=
 \begin{cases}
 f(z)/z,&z\ne0,\\
 f'(0),&z=0
 \end{cases}
\]
在整个 \(\mathbb D\) 全纯：原点附近的级数为 \(f'(0)+f''(0)z/2!+\cdots\)，与去心圆盘上的商相同。固定 \(0<r<1\)，在 \(|z|=r\) 上有 \(|g(z)|\le1/r\)。\(g\) 在闭圆盘 \(\overline{D(0,r)}\) 的邻域全纯，故最大模原理给出
\[
 |g(z)|\le\frac1r\quad(|z|\le r).
\]
现在固定 \(z\in\mathbb D\)，在 \(|z|<r<1\) 中令 \(r\uparrow1\)，得到 \(|g(z)|\le1\)。在 \(z\ne0\) 处乘以 \(|z|\) 得第一项估计，在原点用 \(g(0)=f'(0)\) 得第二项。

若有非零点满足 \(|f(z_0)|=|z_0|\)，则 \(|g(z_0)|=1\)；若 \(|f'(0)|=1\)，则 \(|g(0)|=1\)。两种情形都使 \(g\) 在内部取到其最大模一，所以 \(g\) 是单位模常数 \(c\)，即 \(f(z)=cz\)。反过来，这些函数把 \(\mathbb D\) 映到自身，且两个估计均取等。
\end{proof}

等号条件中的非零点必须保留，因为 \(f(0)=0\) 已使原点的模估计对每个候选函数都取等。如果 \(f\) 不是旋转，证明还给出 \(|f'(0)|<1\) 和 \(|f(z)|<|z|\) 对所有非零 \(z\in\mathbb D\) 成立。不过严格不等号并不自动给出一个在整个圆盘内有效的比例常数 \(q<1\)。例如
\[
 f(z)=\frac{z+z^2}{2}
\]
满足 \(|f(z)|\le |z|\cdot(1+|z|)/2<|z|\) 于 \(0<|z|<1\)，且 \(f'(0)=1/2\)。但当实数 \(r\uparrow1\) 时，\(|f(r)|/r=(1+r)/2\to1\)，所以不存在统一的 \(q<1\) 使整个圆盘上 \(|f(z)|\le q\cdot|z|\)。

\ProseParagraphBreak

中心条件同样有实质作用。常函数 \(f\equiv1/2\) 把单位圆盘映入自身，却不满足 \(|f(z)|\le|z|\) 于原点附近。要处理不固定中心的映射，必须先选取合适的源、目标坐标，把指定点送到原点。第六章将用圆盘自同构完成这一步，得到两点之间的估计。

\ProseParagraphBreak

对于圆心处的比较，仅用平移和缩放即可恢复尺度。设 \(a,b\in\mathbb C\)、\(R,S>0\)，\(F:D(a,R)\to D(b,S)\) 全纯且 \(F(a)=b\)。令
\[
 h(\zeta)=\frac{F(a+R\zeta)-b}{S}\quad(\zeta\in\mathbb D).
\]
它满足 Schwarz 引理的全部条件，故
\[
 |F(z)-b|\le\frac SR\cdot|z-a|,
 \quad |F'(a)|\le\frac SR.
\]
若其中的点值估计在 \(z\ne a\) 取等，或导数估计取等，则 \(h(\zeta)=c\zeta\)、\(|c|=1\)，换回原变量得到
\[
 F(z)=b+c\frac SR(z-a).
\]
因此取等映射正是把两个圆盘作旋转和比例伸缩的仿射映射。

\ProseParagraphBreak

Schwarz 引理还能确定固定圆心的全纯双射。设 \(F:\mathbb D\to\mathbb D\) 全纯且双射，\(F(0)=0\)。由\cref{b3:cor:01-local-injectivity}，\(F'\) 在每点非零，局部逆函数组成全局全纯逆映射 \(H=F^{-1}\)；且 \(H(0)=0\)。分别对 \(F,H\) 使用 Schwarz 引理，有
\[
 |F(z)|\le|z|=|H(F(z))|\le|F(z)|.
\]
所以每个非零点都取等，\(F(z)=cz\)、\(|c|=1\)。归一化后只剩旋转的自由，这将用于共形映射的唯一性；任意圆盘自同构的完整表达式将在后文建立。
